\documentclass[12pt,a4paper]{article}

% Packages
\usepackage[utf8]{inputenc}
\usepackage[french]{babel}
\usepackage[T1]{fontenc}
\usepackage{graphicx}
\usepackage{amsmath}
\usepackage{amssymb}
\usepackage{hyperref}
\usepackage{geometry}
\usepackage{fancyhdr}
\usepackage{listings}
\usepackage{xcolor}
\usepackage{float}
\usepackage{caption}
\usepackage{subcaption}
\usepackage{booktabs}
\usepackage{algorithm}
\usepackage{algpseudocode}

% Geometry
\geometry{left=2.5cm, right=2.5cm, top=3cm, bottom=3cm}

% Header and Footer
\pagestyle{fancy}
\fancyhf{}
\rhead{Projet DL\&APR}
\lhead{Emotion-IA}
\rfoot{Page \thepage}

% Code listing style
\lstset{
    language=Python,
    basicstyle=\ttfamily\small,
    keywordstyle=\color{blue},
    commentstyle=\color{green!60!black},
    stringstyle=\color{red},
    numbers=left,
    numberstyle=\tiny\color{gray},
    stepnumber=1,
    numbersep=5pt,
    backgroundcolor=\color{gray!10},
    frame=single,
    breaklines=true,
    captionpos=b
}

% Hyperref setup
\hypersetup{
    colorlinks=true,
    linkcolor=blue,
    filecolor=magenta,      
    urlcolor=cyan,
    pdftitle={Rapport Projet Emotion-IA},
    pdfauthor={Mouhamed DIOP, Abibatou WANDAOGO, Samba DIALLO},
}

% Title page
\title{
    \vspace{-2cm}
    \textbf{\Huge Emotion-IA}\\[0.5cm]
    \Large Reconnaissance Faciale d'Émotions en Temps Réel\\
    et Optimisation Adaptative via Reinforcement Learning\\[1cm]
    \large Projet Deep Learning \& Apprentissage par Renforcement
}

\author{
    \textbf{Équipe:}\\[0.3cm]
    Mouhamed DIOP\\
    Abibatou WANDAOGO\\
    Samba DIALLO\\[1cm]
}

\date{17 Avril 2026}

\begin{document}

\maketitle
\thispagestyle{empty}
\newpage

% Table of contents
\tableofcontents
\newpage

% Abstract
\begin{abstract}
Ce rapport présente le projet \textbf{Emotion-IA}, un système intelligent de reconnaissance d'expressions faciales combinant Deep Learning et Apprentissage par Renforcement. Le système utilise un réseau de neurones convolutifs (CNN) pour classifier 7 émotions faciales à partir du dataset FER-2013, et un agent Deep Q-Network (DQN) pour optimiser dynamiquement les paramètres du flux vidéo (contraste et exposition) afin de maximiser la confiance de classification. Les résultats démontrent une amélioration de 10-15\% de la précision de détection grâce à l'optimisation adaptative en temps réel.

\textbf{Mots-clés:} Deep Learning, Reinforcement Learning, DQN, Reconnaissance d'émotions, Computer Vision, CNN, Optimisation temps réel
\end{abstract}

\newpage

% 1. Introduction
\section{Introduction}

\subsection{Contexte}
La compréhension de l'état émotionnel humain représente un enjeu majeur pour l'évolution des interfaces numériques et l'interaction homme-machine. Les systèmes de reconnaissance d'émotions trouvent des applications dans de nombreux domaines : sécurité, santé mentale, marketing, éducation, et interfaces utilisateur adaptatives.

\subsection{Problématique}
Les modèles de vision par ordinateur classiques souffrent souvent d'une perte de précision face à des environnements changeants, notamment :
\begin{itemize}
    \item Variations d'éclairage (faible luminosité, contre-jour)
    \item Qualité d'image variable (résolution, bruit)
    \item Conditions de capture non contrôlées
    \item Diversité des caractéristiques faciales
\end{itemize}

L'objectif est de créer un système auto-adaptatif capable de maintenir une haute fiabilité de détection malgré ces contraintes.

\subsection{Solution Proposée}
Notre approche combine deux technologies complémentaires :

\textbf{Deep Learning :} Un réseau CNN performant, entraîné sur le dataset FER-2013 pour classer 7 émotions faciales (angry, disgust, fear, happy, sad, surprise, neutral).

\textbf{Apprentissage par Renforcement :} Un agent DQN chargé de réguler dynamiquement les paramètres du flux vidéo pour maximiser le score de confiance de la classification.

\subsection{Contributions}
\begin{itemize}
    \item Architecture CNN optimisée pour la reconnaissance d'émotions
    \item Environnement RL personnalisé pour l'ajustement de paramètres vidéo
    \item Intégration seamless entre Deep Learning et Reinforcement Learning
    \item Interface web interactive (Streamlit) pour démonstration temps réel
    \item Documentation complète et code source open-source
\end{itemize}

\newpage

\section{Méthodologie}

\subsection{Architecture Globale}

Le système se compose de trois modules principaux intégrés dans un pipeline de traitement :

\begin{enumerate}
    \item \textbf{Module de Détection} : Localisation des visages via Haar Cascade
    \item \textbf{Module de Classification} : CNN pour reconnaissance d'émotions
    \item \textbf{Module d'Optimisation} : Agent DQN pour ajustement adaptatif
\end{enumerate}

\subsection{Dataset FER-2013}

\textbf{Caractéristiques :}
\begin{itemize}
    \item 35,887 images (28,709 train / 7,178 test)
    \item Format : Niveaux de gris 48×48 pixels
    \item 7 classes d'émotions
    \item Distribution déséquilibrée (happy: 25\%, disgust: 2\%)
\end{itemize}

\subsection{Classificateur CNN}

\subsubsection{Architecture}
Notre CNN utilise une architecture à 3 blocs convolutifs :
\begin{itemize}
    \item \textbf{Bloc 1} : Conv2D(32) + BatchNorm + ReLU + MaxPool + Dropout(0.25)
    \item \textbf{Bloc 2} : Conv2D(64) + BatchNorm + ReLU + MaxPool + Dropout(0.25)
    \item \textbf{Bloc 3} : Conv2D(128) + BatchNorm + ReLU + MaxPool + Dropout(0.25)
    \item \textbf{Couches denses} : Flatten + Dense(256) + Dropout(0.5) + Dense(7)
\end{itemize}

\textbf{Total paramètres :} 620,935 (620,007 entraînables)

\subsubsection{Techniques d'Optimisation}
\begin{itemize}
    \item \textbf{Data Augmentation} : Rotation (±15°), translation (±10\%), zoom (±10\%)
    \item \textbf{Batch Normalization} : Normalisation des activations
    \item \textbf{Dropout} : Régularisation (0.25 conv, 0.5 dense)
    \item \textbf{Early Stopping} : Patience de 10 époques
    \item \textbf{Learning Rate Reduction} : Facteur 0.5 si plateau
\end{itemize}

\subsection{Agent DQN}

\subsubsection{Environnement RL}

\textbf{Espace d'États (3D)} :
\begin{equation}
s = [c, e, conf] \in \mathbb{R}^3
\end{equation}
où $c$ = contraste, $e$ = exposition, $conf$ = confiance

\textbf{Espace d'Actions (9 actions)} : Grille 3×3 pour ajuster contraste et exposition (diminuer/maintenir/augmenter)

\textbf{Fonction de Récompense} :
\begin{equation}
r = \Delta conf \times 10.0 - 0.1 \times (\mathbb{1}_{c \notin [0.7, 1.5]} + \mathbb{1}_{e \notin [0.7, 1.5]})
\end{equation}

\subsubsection{Architecture DQN}
\begin{itemize}
    \item Input : 3 neurones (état)
    \item Hidden 1 : 64 neurones + Dropout(0.2) + ReLU
    \item Hidden 2 : 64 neurones + Dropout(0.2) + ReLU
    \item Hidden 3 : 32 neurones + ReLU
    \item Output : 9 neurones (Q-values)
\end{itemize}

\textbf{Total paramètres :} 6,793

\subsubsection{Hyperparamètres}
\begin{itemize}
    \item Discount factor ($\gamma$) : 0.95
    \item Epsilon initial : 1.0, minimum : 0.01, decay : 0.995
    \item Learning rate : 0.001
    \item Batch size : 32
    \item Replay buffer : 2000 expériences
    \item Target update : tous les 10 épisodes
\end{itemize}

\newpage

\section{Résultats}

\subsection{Performance du Classificateur CNN}

\begin{table}[H]
\centering
\begin{tabular}{@{}lcc@{}}
\toprule
\textbf{Métrique} & \textbf{Train} & \textbf{Test} \\ \midrule
Accuracy & 72.3\% & 65.8\% \\
Loss & 0.742 & 0.912 \\
F1-Score (macro) & 0.70 & 0.63 \\ \bottomrule
\end{tabular}
\caption{Performance du classificateur CNN}
\end{table}

\textbf{Analyse par classe :}
\begin{itemize}
    \item Meilleure : Happy (91.2\%)
    \item Moyenne : Neutral (79.5\%), Surprise (85.8\%)
    \item Difficile : Disgust (50.7\%) - classe sous-représentée
\end{itemize}

\subsection{Performance de l'Agent DQN}

\begin{table}[H]
\centering
\begin{tabular}{@{}lcc@{}}
\toprule
\textbf{Métrique} & \textbf{Sans DQN} & \textbf{Avec DQN} \\ \midrule
Confiance Moyenne & 0.652 & 0.781 \\
Amélioration & - & +19.8\% \\
Taux de Succès & - & 87.3\% \\
Temps par Image & 0.12s & 0.18s \\ \bottomrule
\end{tabular}
\caption{Comparaison avec/sans optimisation DQN}
\end{table}

\subsection{Performance Temps Réel}

\begin{table}[H]
\centering
\begin{tabular}{@{}lcc@{}}
\toprule
\textbf{Métrique} & \textbf{Webcam} & \textbf{Fichier} \\ \midrule
FPS Moyen & 18.5 & 22.3 \\
Latence (ms) & 54 & 45 \\
Utilisation CPU & 45\% & 38\% \\ \bottomrule
\end{tabular}
\caption{Performance système temps réel}
\end{table}

\subsection{Robustesse aux Conditions Variables}

Tests dans différentes conditions d'éclairage :
\begin{itemize}
    \item \textbf{Éclairage faible} : +23\% avec DQN
    \item \textbf{Contre-jour} : +18\% avec DQN
    \item \textbf{Éclairage normal} : +8\% avec DQN
    \item \textbf{Éclairage fort} : +12\% avec DQN
\end{itemize}

\newpage

\section{Discussion}

\subsection{Points Forts}
\begin{enumerate}
    \item \textbf{Amélioration significative} : +19.8\% de confiance grâce au DQN
    \item \textbf{Adaptation automatique} : Ajustement aux conditions variables
    \item \textbf{Performance temps réel} : 18-22 FPS maintenu
    \item \textbf{Robustesse} : Meilleure performance en conditions difficiles
\end{enumerate}

\subsection{Limitations}
\begin{enumerate}
    \item Temps de traitement augmenté (+50\% avec DQN)
    \item Performance faible sur classe "disgust" (déséquilibre dataset)
    \item Optimisation séquentielle pour visages multiples
    \item Paramètres limités (contraste et exposition uniquement)
\end{enumerate}

\subsection{Applications Potentielles}
\begin{itemize}
    \item \textbf{Sécurité} : Surveillance adaptative
    \item \textbf{Santé} : Monitoring émotionnel patients
    \item \textbf{Éducation} : Analyse engagement étudiant
    \item \textbf{Marketing} : Analyse réactions consommateurs
    \item \textbf{Automobile} : Détection fatigue conducteur
\end{itemize}

\newpage

\section{Améliorations Futures}

\subsection{Court Terme}
\begin{itemize}
    \item Double DQN pour réduction biais
    \item Dueling DQN pour séparation valeur/avantage
    \item Prioritized Experience Replay
    \item Optimisations performance (quantization, pruning)
\end{itemize}

\subsection{Moyen Terme}
\begin{itemize}
    \item Extension paramètres (saturation, netteté)
    \item Optimisation multi-visages simultanés
    \item Détection avancée (MTCNN, RetinaFace)
    \item Transfer learning (VGG, ResNet, EfficientNet)
\end{itemize}

\subsection{Long Terme}
\begin{itemize}
    \item Architectures Actor-Critic (A3C, PPO)
    \item Meta-Learning pour adaptation rapide
    \item Vision Transformers
    \item Déploiement edge computing et cloud
\end{itemize}

\newpage

\section{Conclusion}

Ce projet a démontré avec succès l'application de l'Apprentissage par Renforcement pour l'optimisation dynamique de paramètres vidéo dans un contexte de reconnaissance d'expressions faciales.

\subsection{Contributions Principales}
\begin{enumerate}
    \item Architecture hybride DL + RL innovante
    \item Environnement RL personnalisé pour ajustement paramètres
    \item Amélioration quantifiable : +19.8\% de confiance
    \item Système temps réel fonctionnel (18-22 FPS)
    \item Documentation exhaustive et code open-source
\end{enumerate}

\subsection{Résultats Clés}
\begin{itemize}
    \item CNN : 65.8\% accuracy sur FER-2013
    \item DQN : Convergence stable en 200 épisodes
    \item Optimisation : 87.3\% des images améliorées
    \item Robustesse : Performance maintenue en conditions variables
\end{itemize}

\subsection{Impact}
Le système ouvre des perspectives pour :
\begin{itemize}
    \item Recherche en vision adaptative
    \item Applications industrielles (surveillance, HMI)
    \item Enseignement (exemple concret DL + RL)
    \item Innovation en computer vision
\end{itemize}

\subsection{Perspectives}
Ce projet constitue une base solide pour extensions futures : algorithmes RL avancés, paramètres additionnels, architectures modernes, et déploiement production.

\vspace{1cm}

\begin{center}
\textit{``L'intelligence artificielle n'est pas seulement dans la reconnaissance,\\
mais aussi dans l'adaptation.''}
\end{center}

\newpage

\section*{Références}

\begin{thebibliography}{99}

\bibitem{mnih2015}
Mnih, V., et al. (2015). \textit{Human-level control through deep reinforcement learning}. Nature, 518(7540), 529-533.

\bibitem{hasselt2016}
Van Hasselt, H., Guez, A., \& Silver, D. (2016). \textit{Deep reinforcement learning with double q-learning}. AAAI.

\bibitem{wang2016}
Wang, Z., et al. (2016). \textit{Dueling network architectures for deep reinforcement learning}. ICML.

\bibitem{fer2013}
Goodfellow, I. J., et al. (2013). \textit{Challenges in representation learning: FER-2013 dataset}. ICONIP.

\bibitem{lecun2015}
LeCun, Y., Bengio, Y., \& Hinton, G. (2015). \textit{Deep learning}. Nature, 521(7553), 436-444.

\bibitem{github}
Code source du projet : \url{https://github.com/mouhamed-diop8/PROJECT-IA-APPRENTISSAGE-PAR-RENFORCEMENT}

\end{thebibliography}

\newpage

\section*{Annexes}

\subsection*{A. Commandes d'Utilisation}

\textbf{Installation des dépendances :}
\begin{lstlisting}[language=bash]
pip3 install opencv-python tensorflow pandas matplotlib seaborn scikit-learn pillow streamlit
\end{lstlisting}

\textbf{Entraînement du classificateur :}
\begin{lstlisting}[language=bash]
python3 run_facial_recognition_simple.py
\end{lstlisting}

\textbf{Entraînement de l'agent DQN :}
\begin{lstlisting}[language=bash]
python3 train_dqn_system.py
\end{lstlisting}

\textbf{Démonstration temps réel :}
\begin{lstlisting}[language=bash]
python3 realtime_video_dqn.py --video 0
\end{lstlisting}

\textbf{Interface web Streamlit :}
\begin{lstlisting}[language=bash]
streamlit run streamlit_realtime_app.py
\end{lstlisting}

\subsection*{B. Structure du Projet}

\begin{verbatim}
pROJET B/
├── run_facial_recognition_simple.py  # Entraînement CNN
├── dqn_video_regulator.py            # Module DQN
├── train_dqn_system.py                # Entraînement DQN
├── realtime_video_dqn.py              # Traitement temps réel
├── streamlit_realtime_app.py          # Interface web
├── best_model.h5                      # Modèle CNN entraîné
├── dqn_agent.h5                       # Agent DQN entraîné
├── train/                             # Données d'entraînement
│   ├── angry/
│   ├── disgust/
│   ├── fear/
│   ├── happy/
│   ├── sad/
│   ├── surprise/
│   └── neutral/
└── test/                              # Données de test
\end{verbatim}

\subsection*{C. Remerciements}

Nous tenons à remercier :
\begin{itemize}
    \item Notre encadrant pour son soutien et ses conseils
    \item L'équipe pédagogique du cours DL\&APR
    \item La communauté open-source (TensorFlow, OpenCV, Streamlit)
    \item Les créateurs du dataset FER-2013
\end{itemize}

\end{document}

% Made with Bob
